\documentclass{article}
\usepackage[utf8]{inputenc}
\usepackage[T1]{fontenc}
\usepackage[margin=2cm]{geometry}
\usepackage{hyperref}
\sloppy
\begin{document}
\section*{Historical source}
Encyclopaedia Britannica scan, source row 1344. Parse 3.5 extracted the text. Super 3.5 VL corrected it from an enlarged crop of each block. It can still contain recognition errors.
Page headers, footers, and page numbers are removed.

\subsection*{Untitled continuation (E0116)}
\textit{Continues an entry from an earlier page.}

H O R [ 664 ] H O R

Horse.

be totally closed up, or the cranium in that part not be
so hard as the rest for some years after. This open-
ness is found to be increased upon the child's catching
cold. When the disease continues long, it is reput-
ed a sign of weakness and short life. In this case,
it is usual to rub the head now and then with warm
rum or brandy, mixed with the white of an egg and
palm-oil. Sometimes the disorder arises from a col-
lection of waters in the head called an hydrocephalus.

lection of waters in the head called an hydrocephalus.
Stone-Horse. See Stallion.
Horse Tail. See Equisetum and Equus.

Stone-Horse. See Stallion.
Horse Tail. See Equisetum and Ephedra.
Mare's Tail. See Hippuris.

Horse Tail. See Equisetum and Ephedra.
Horse-Vetch. See Hippocrepis.
War Horse. The proper rules for shoeing a

War-Horse. The proper rules for choosing a horse for service in war, are these: he should be tall in stature, with a comely head, and out-swelling forehead. His eye should be bright and sparkling, and the white part of it covered by the eye brow. The ears should be small, thin, short, and pricking; or if long, they should be moveable with ease, and well carried. The neck should be deep, and the breath large and swelling. The ribs bending, the chine broad and straight, and the buttocks round and full. The tail should be high and broad, neither too thick nor too thin; the thigh swelling; the leg broad and flat, and the pastern short. When such a horse is chosen, he must be kept high during the time of his teaching, that he may be full of vigour. His food must be sweet hay, and good clean oats, or two parts of oats and one part of beans or pease, well dried and hardened. The quantity should be half a peck in the morning, and the same quantity at noon and in the evening. Upon his resting days he is to be drest between five and six in the morning, and watered at seven or eight. In the evening he is to be drest at four, and watered about five, and he must always have provender given him after watering; he must be littered about eight, and then must have food given him for all night. The night before he is ridden all his hay is to be taken away about nine o'clock, and he must have a handful or two of oats about four in the morning; when he has eaten these, he is to be turned upon the straw, and rubbed very well with dry cloths; then saddled, and made fit for his exercise. When he has performed this, he is to be brought sweating into the stable, and rubbed down with dry whips. When this has been done, the saddle is to be taken off, and he is to be rubbed down with dry cloths; the housing cloth is then to be laid on; and the saddle being again laid on, he is to be walked gently about till thoroughly cool. After this, he must stand without meat two or three hours, then he must be fed; and in the afternoon he is to be rubbed and drest as before, and watered in the usual manner.

Horse Worm, in natural history, a species of fly-
worm called also bot, produced of eggs deposited by a
two-winged fly of the shape and size of the humble
bee in the intestines of horses. See Borrs.

River-Horse, in zoology. See HIPPOPOTAMUS.
Horse is also used in the military language, to ex-

Horse is also used in the military language, to ex-
press the cavalry; or the body of soldiers who serve on
horseback.

The horse includes horse guards, horse grenadiers,
and troopers. Dragoons are also frequently compre-
hended under this name, though they fight on foot:
of these there are now 18 regiments ; besides three re-

ments of dragoon-guards raised in 1685. See Gra-
nadier, Dragoons, and Guards.

\subsection*{ADIER, DRAGOONS, and GUARDS (E0117)}
ADIER, DRAGOONS, and GUARDS.
Master of the Horse. See MASTER.
Light Horse, are regiments of cavalry

Light-Horse, are regiments of cavalry, mounted on
light swift horses, whose men are small and lightly ac-
coutred. They were first raised in 1757. The deno-
mination arose hence, that anciently they were lightly
armed, in comparison of the royal guards, which were
armed at all points.

Hungarian Horse. See Hussar.

Horse is also a term used in various arts and manufactories, for something that helps to sustain their work from the ground, for the more commodious working at it.

The horse used by tanners and fkinners, also called
the leg, is a piece of wood cut hollow and roundish,
four or five feet long, and placed alope; upon which
they pare their fkins to get off the dirt, hair, flesh, \&c.

Horse is also used in carpentry, for a piece of wood
joined across two other perpendicular ones, to sustain
the boards, planks, \&c. which make bridges over small
rivers; and on divers other occasions.

Horse, in sea language, is the name of a rope reaching from the middle of a yard to its extremity, or what is called the yard-arm, and depending about two or three feet under the yard, for the sailors to tread upon whilst they are loosing, reefing, or furling the sails, rigging out the studding-sail booms, \&c. In order, therefore, to keep the horse more parallel to the yard, it is usually suspended to it at proper distances, by certain ropes called stirrups, which hang about two feet under the yard, having an eye in their lower ends through which the horse passes.

Horse is also a thick rope, extended in a perpendicular direction near the fore or after-side of a mast, for the purpose of hoisting or extending some fail upon it. When it is fixed before a mast, it is calculated for the use of a fail called the square-fail, whose yard being attached to the horse, by means of a traveller or bull's eye, which slides up and down occasionally, is retained in a steady position; either when the fail is set, or whilst it is hoisting or lowering. When the horse is placed abaft or behind a mast, it is intended for the try-fail of a snow, and is accordingly very rarely fixed in this position, except in those fleets of war which occasionally assume the form of snow, in order to deceive the enemy.

Horse is also a cant name introduced into the management of lotteries, for the chance or benefit of a ticket or number for one or more days, upon condition, if it be drawn a prize within the time covenanted for, of returning to the seller an undrawn ticket.—To determine the value of a horse; multiply the amount of the prizes in the lottery by the time the horse is hired for; and from the product subtract the amount of the number of prizes by the value of an undrawn ticket into the time of the horse; the remainder being divided by the number of tickets into the whole time of drawing, the quotient is the value of the horse. See Lottery.

\_Horas-Bread.\_ See Barxad.

Horse-Bread. See BREAD.
Horse-Dung, in gardening, is of great use in ma-
king hot-beds, for the raising all sorts of early crops;
as fallading, cucumbers, melons, asparagus, \&c. for
which


\section*{Source record}
\begin{description}
\item[Scan ID] \texttt{britannica:309f765b1ecc83e0a7f7762e44b4a927a7555443:train:1344}
\item[Image file] \texttt{page-1344-b5c2505d4784.jpg}
\item[Image SHA-256] \texttt{b5c2505d4784ce9e909dbebeedb4858735d3550cf3e7506a2a92c391393e8cc2}
\item[Dataset] NationalLibraryOfScotland/encyclopaedia\_britannica\_illustrated at revision \texttt{309f765b1ecc83e0a7f7762e44b4a927a7555443} (CC0-1.0)
\item[Extraction] model \texttt{nvidia/nemotron-3.5-parse}, raw reply \texttt{extract/raw/page-1344-b5c2505d4784.json}
\item[Raw reply SHA-256] \texttt{2a387187c1222343478b69608413f0bc7ccf71fe57fc6717a79fbcfd801a399f}
\item[Correction] model \texttt{nvidia/Nemotron-3.5-Super-VL}, block crops, raw replies \texttt{correct/raw/crop\_nothink/page-1344-b5c2505d4784-b*.json}; parse\_fallback: 1, super\_vl: 22
\item[Running header] none
\end{description}
\end{document}
